\section{Erd\H{o}s Problem \#373: independent continuation}
\noindent Date: 2026-09-23. Research attempt; no claim of a new published result.
\subsection{1) TARGET AND SOURCE AUDIT}
Are there only finitely many solutions of $n!=a_1!\cdots a_k!$ with $n-1>a_1\ge\cdots\ge a_k\ge2$?

All supplied versions considered: \texttt{373.tex}.
The source forces $a_1$ close to $n$ through a prime-gap theorem. I derive complementary restrictions on the remaining factorials using elementary size and valuation arguments; no stronger prime-gap result is assumed.
\subsection{2) WORK}
Put $h=n-a_1\ge2$. Cancelling $a_1!$ gives
\[
\prod_{i=2}^k a_i!=(a_1+1)\cdots n\le n^h.
\]
Each factor on the left is at least two, so
\[
k-1\le h\log_2 n. \tag{1}
\]
Also for $a_2\ge2$ the last $\lfloor a_2/2\rfloor$ factors in $a_2!$ are all at least $a_2/2$, giving
\[
\lfloor a_2/2\rfloor\log(a_2/2)\le h\log n. \tag{2}
\]
For any fixed $h$, inequality (2) implies $a_2=O_h(\log n/\log\log n)$: inserting $a_2\ge C h\log n/\log\log n$ makes its left side greater than $h\log n$ for a sufficiently large absolute $C$ and large $n$.

There is also an exact prime restriction. If $p$ is prime with $n/2<p\le a_1$, then $v_p(n!)=v_p(a_1!)=1$. Every other factorial must have $p$-valuation zero, so $a_2<p$. If $a_1<p\le n$ instead, the equation is impossible. Thus a hypothetical solution with a fixed short terminal block must express that block's product using factorials whose arguments are only logarithmic in $n$, despite the block entries being of order $n$.
\subsection{3) ADVERSARIAL VERIFICATION}
Repeated factorial arguments are permitted in the problem and in (1). The bound for fixed $h$ is not uniform in a growing $h$ unless its dependence is retained; the source only gives a growing upper bound on $h$. Smooth products of large consecutive integers are not excluded merely by these inequalities.
\subsection{4) FINAL}
\textbf{UNRESOLVED.}
No contradiction for all large $n$ is obtained, even when $h$ is fixed. The remaining issue is the arithmetic feasibility of the residual factorial product, not just the closeness of $a_1$ to $n$.
\subsection{5) SOURCES CHECKED}
No external theorem is used in the new argument. The supplied source was read in full; the saved Tao status snapshot is dated 2026-09-22.
\subsection{6) COMPLETION ESTIMATE}
This is a conservative subjective measure of progress toward the original question, not a probability of correctness or a claim that the remaining work is routine.
\textbf{COMPLETION: 10\%}
